\documentclass[11pt]{article}
\usepackage[a4paper,margin=27mm]{geometry}
\usepackage[T1]{fontenc}
\usepackage[utf8]{inputenc}
\usepackage{lmodern,amsmath,amssymb,amsthm,mathrsfs,booktabs,array,microtype}
\usepackage[colorlinks=true,linkcolor=blue,citecolor=blue,urlcolor=blue]{hyperref}
\hypersetup{pdftitle={Fixed-branch swim distance for triangular color-code matching},
pdfauthor={Color-code soft-output research project}}
\theoremstyle{definition}
\newtheorem{definition}{Definition}[section]
\theoremstyle{plain}
\newtheorem{lemma}[definition]{Lemma}
\newtheorem{proposition}[definition]{Proposition}
\newtheorem{theorem}[definition]{Theorem}
\newtheorem{corollary}[definition]{Corollary}
\theoremstyle{remark}
\newtheorem{remark}[definition]{Remark}
\newtheorem{openissue}[definition]{Open issue}
\newcommand{\F}{\mathbb F_2}
\newcommand{\C}{\mathcal C}
\newcommand{\LL}{\mathcal L}
\newcommand{\bd}{v_{\mathrm{bdry}}}
\newcommand{\eps}{\epsilon}
\newcommand{\pred}{\mathrm{pred}}
\newcommand{\dang}{\mathrm{dang}}
\newcommand{\missF}{\mathrm{missF}}
\newcommand{\missE}{\mathrm{missE}}
\newcommand{\MWPM}{\operatorname{MWPM}}
\title{Fixed-branch swim distance for triangular color-code matching}
\author{Color-code soft-output research note}
\date{16 September 2026}
\begin{document}
\maketitle
\begin{abstract}
We specify the physical and matching-lattice conventions for perfect-syndrome
decoding of standard triangular 6.6.6 color codes. For every odd distance
$d\geq3$ and fixed color $c$, we classify all physical qubits whose
monochromatic matching edges meet the ordinary artificial boundary node.
They form the $d$ qubits on the closed $c$-colored side and the single
opposite corner. The two sets are distinguished by a missing $c$ face and a
missing $c$ edge, respectively. An explicit coordinate proof covers the
entire family, including every corner. We then resolve these types
into two virtual vertices, prove connectivity and an explicit quotient
back to the ordinary graph, and establish compatibility with matching
when both terminals are free. A linear qubit-support map is defined on
all resolved edge chains, with an exact factorization of the physical
syndrome through graph incidence. Physical logical-boundary overlap then
proves that terminal-joining relative chains are precisely the nontrivial
logical class and closed chains are generated exactly by face stabilizers of the other
two colors. The resulting induced binary complex realizes the
physical logical quotient, establishing the fixed-branch topology of a
Meister-type modified decoding graph. A free-boundary optimal dual
specifies decoder clusters, whose exact interval contraction gives a
computable per-color swim distance. We prove its physical logical
interpretation and a certified stage-2 minimum-representative gap bound.
The result does not identify a full-decoder gap or a logical-class LLR.
The circuit-level part starts from detector and observable maps of the
retained decomposed model. It proves a fixed-fiber quotient, a canonical
logical-cover algorithm, a gated boundary cut, compatible labelled
contraction, and a certified representative bound. A controlled
perfect-measurement limit recovers the preceding theory exactly.
\end{abstract}

\part{Perfect-measurement fixed-fiber theory}
\section{Scope and physical code}\label{sec:physical}
All supports and incidence maps are binary. Addition of supports means
symmetric difference. We analyze $X$ errors using measured $Z$ checks.
Exchanging $X$ and $Z$ gives the other CSS sector; this does not assert
statistical independence for arbitrary Pauli noise. Measurements are
perfect in the first part (Sections~1--13); Part~\ref{part:circuit}
specifies the circuit-level extension. Holes, twists, Pauli-type boundaries, other tilings such as
4.8.8, and the degenerate distance-one patch are excluded.

\begin{definition}[Physical lattice and colors]\label{def:physical}
Let $\LL_{2D}$ be the standard triangular patch of the hexagonal color-code
lattice, with
\[
 Q=\Delta_0(\LL_{2D}),\quad E=\Delta_1(\LL_{2D}),\quad
 F=\Delta_2(\LL_{2D}),\quad \C=\{r,g,b\}.
\]
These are physical data vertices, primal edges, and real faces.
Qubits occupy $Q$, not faces or edges. Write
$F_c=\Delta_2^{(c)}(\LL_{2D})$ and $E_c=\Delta_1^{(c)}(\LL_{2D})$.
The bulk lattice is trivalent with properly three-colored faces; the
finite patch has degree-two corners.

The three physical boundary arcs are $\Gamma_r,\Gamma_g,\Gamma_b$.
A boundary of color $a$ is adjacent only to real faces of the other colors.
An interior primal edge separating face colors $a,b$ has the third color.
A boundary primal edge on $\Gamma_a$ with incident real face of color $b$
also has the third color. In this rule the absent exterior sector is
labelled by $a$; it is not a physical check.
\end{definition}

These are the physical and boundary conventions of
Lee--Li--Bartlett \cite[Sec.~2.1, pp.~3--4]{Lee}.
Section~\ref{sec:family} fixes the precise finite-family realization.
The corner exception matters when calling the lattice trivalent.

\begin{definition}[Checks, syndrome, and physical logical class]\label{def:checks}
For $a\in A=\F^Q$ and $f\in F$, let
\[
 P_X(a)=\prod_{q\in Q}X_q^{a_q},\quad
 P_Z(a)=\prod_{q\in Q}Z_q^{a_q},\qquad
 S_f^X=\prod_{q\in f}X_q,\quad S_f^Z=\prod_{q\in f}Z_q .
\]
Here $q\in f$ means face-support incidence. Let $U=\F^F$ and define
\[
 H:A\longrightarrow U,\qquad (Ha)_f=\sum_{q\in f}a_q,\qquad
 J=H^{\mathsf T}:U\longrightarrow A.
\]
The stabilizer group $\mathcal S$ is generated by the real-face
$S_f^X,S_f^Z$. Faces have even weight, and distinct intersecting faces
share two qubits, so $HJ=0$. The code space is their common $+1$
eigenspace. An $X$ error of support $a$ has syndrome $\sigma_Z=Ha$:
the outcome of $S_f^Z$ is $(-1)^{(Ha)_f}$. We also write
$\sigma_Z\subseteq F$ for the syndrome support and
$\sigma_Z^{(c)}=\sigma_Z\cap F_c$.
\end{definition}

The checks are those of \cite[Eq.~(1)]{Lee} and
\cite[Eq.~(1)]{Bombin}. The physical CSS complex is
\[
 U\xrightarrow{J}A\xrightarrow{H}U,\qquad HJ=0.
\]
Thus $S_X=\operatorname{im}J$, $Z_X=\ker H$, and
$\LL_X=Z_X/S_X$ describe physical stabilizer supports, zero-syndrome
supports, and logical classes, respectively. They are not the homology
of an unspecified matching graph. Same-syndrome corrections $a,a'$
are equivalent exactly when $a+a'\in S_X$. Syndrome validity alone gives
$a+a'\in Z_X$.

The standard triangular code encodes one qubit, so $\dim\LL_X=1$
\cite[triangular-code construction, pp.~2--3]{Bombin}.
Logical $X$ and $Z$ representatives can each be supported on a complete
physical side \cite[Sec.~2.1, p.~4]{Lee}.
Write $Q_a=Q\cap\Gamma_a$, including both corners of that side.
Lemma~\ref{lem:incidence} verifies $|Q_a|=d$.
The same-side representatives $P_X(Q_a),P_Z(Q_a)$ anticommute since
$d$ is odd. Fixing a nontrivial logical-$Z$ support $\ell_Z$ gives
the physical logical parity $\lambda_X(z)=\ell_Z^{\mathsf T}z$ on
$Z_X$. No equality with a matching-path endpoint parity is assumed.

\section{Colored strings and boundary termination}\label{sec:strings}
This section concerns physical strings, not logical paths of $\LL_c^*$.
For a primal edge $e=\{u,v\}\in E_c$, the operator $X_uX_v$ has syndrome
$H(u+v)$, with $u,v$ also denoting basis vectors of $A$.
Every non-$c$ face contains either both endpoints or neither, so those
contributions cancel. Remaining excitations are at the $c$ faces at the
ends; a boundary may omit such a face. This local incidence property
also follows from Lemma~\ref{lem:incidence}.

A colored Pauli string transports a charge by composing these elementary
two-qubit edge operators along a path in the usual shrunk lattice of that
color. Consecutive segments share a $c$-face endpoint, whose syndrome
cancels in the binary sum. The physical support is the symmetric
difference of the endpoint pairs of the traversed primal $c$ edges.
Only uncancelled real endpoints are excited. Exchanging $X$ and $Z$
exchanges the excited check type. These are the colored strings of
\cite[Eq.~(4), pp.~2--3]{Bombin}; their shrunk lattice is not being
identified with either matching lattice defined below.

A boundary of color $c$ condenses the corresponding colored charges:
a $c$ string can terminate there without a measured $c$-check excitation
\cite[boundary discussion, p.~3]{Bombin},
\cite[Secs.~4.2--4.3, Fig.~6 and Table~2]{Kesselring}.
This does not mean that all endpoints at one virtual matching node are
logically equivalent. At a corner where colors $a,b$ meet, the only real
face has the remaining color $c$; a single-qubit $X$ there excites only
that real $Z$ check.

Strings can meet at junctions to form string-nets. The three color charges
of a fixed Pauli type can fuse to vacuum
\cite[Secs.~3.3--3.4]{Kesselring}.
Appropriate three-string nets give triangular-code logical representatives
\cite[Eqs.~(5)--(6), Fig.~3]{Bombin}.
Multiplying face stabilizers changes a support by an element of
$\operatorname{im}J$ and realizes stabilizer deformations. A logical
operator is therefore an equivalence class, not a preferred curve.
Neither arbitrary graph cycles nor open stage-2 paths are declared to be
such operators here.

\section{A combinatorial realization of every allowed patch}\label{sec:family}
\begin{definition}[Standard odd-distance family]\label{def:coords}
Let $d=2t+1$, $t\geq1$, and $L=3t$. In triangular-lattice coordinates set
\[
 \Omega_t=\{(i,j)\in\mathbb Z_{\geq0}^2:i+j\leq L\},\quad
 F_t=\{(i,j)\in\Omega_t:i-j\equiv2\pmod3\},\quad
 Q_t=\Omega_t\setminus F_t.
\]
Identify $Q_t$ with physical qubits and each $f\in F_t$ with a real face.
The face color is $\chi(j\bmod3)$, with
$\chi(0)=g,\ \chi(1)=b,\ \chi(2)=r$.
Use the cyclic neighbor list
\[
 N=\big((-1,1),(0,1),(1,0),(1,-1),(0,-1),(-1,0)\big).
\]
The face at $f$ has cyclic support $(f+N)\cap Q_t$, retaining this order.
Consecutive retained qubits, including the wraparound pair, define its
primal edges; shared edges are identified. This produces hexagonal bulk
faces and quadrilateral boundary faces. The sides and corners are
\begin{align*}
 &\Gamma_r:\ j=0,\qquad\Gamma_g:\ i=0,\qquad\Gamma_b:\ i+j=L,\\
 &q_{rg}=(0,0),\qquad q_{rb}=(L,0),\qquad q_{gb}=(0,L).
\end{align*}
A side is the polygonal arc through its retained qubits. Excluded
face-center sites on its supporting line are not qubits.
\end{definition}

This specifies the standard triangular 6.6.6 family used in this note.
Equivalently one starts from the hexagonal tiling with residue-two face
centers and truncates the boundary polygons along the three sides.
The data/check coordinates agree with the author's triangular
implementation \cite{Implementation} under $x=4i+2j,\ y=j$, before the
check-ancilla display displacement. That implementation fixes the
coordinate convention; the following proof is symbolic for every
$t\geq1$, not a finite test of the implementation.

\begin{definition}[Local incidence sets]\label{def:local}
For $q\in Q$ and $c\in\C$, put
\[
 \mathcal F_c(q)=\{f\in F_c:q\in f\},\qquad
 \mathcal E_c(q)=\{e\in E_c:q\in e\},\qquad
\]
\[
 n_c^F(q)=|\mathcal F_c(q)|,\qquad n_c^E(q)=|\mathcal E_c(q)|.
\]
\end{definition}

\begin{lemma}[Complete local incidence]\label{lem:incidence}
For all patches in Definition~\ref{def:coords}, the exhaustive local
incidence table is
\begin{center}
\begin{tabular}{@{}lll@{}}
\toprule
Physical location & Real-face colors & Primal-edge colors\\
\midrule
Bulk & $r,g,b$ & $r,g,b$\\
Open side $\Gamma_a$ & $\C\setminus\{a\}$ & $r,g,b$\\
Corner $q_{ab}$ & $\C\setminus\{a,b\}$ & $a,b$\\
\bottomrule
\end{tabular}
\end{center}
Each listed color occurs once. Furthermore $|Q_a|=d$,
$Q_a\cap Q_b=\{q_{ab}\}$ for $a\ne b$, and no qubit lies on all three sides.
\end{lemma}
\begin{proof}
For $q=(i,j)$ with $i-j\equiv0\pmod3$, its possible face centers are
\[
 (i-1,j),\qquad(i,j+1),\qquad(i+1,j-1).
\]
They are absent exactly when $i=0$, $i+j=L$, and $j=0$, respectively.
For $i-j\equiv1\pmod3$, the possible centers are
\[
 (i+1,j),\qquad(i,j-1),\qquad(i-1,j+1),
\]
absent exactly when $i+j=L$, $j=0$, and $i=0$, respectively.
Their row residues are all different, hence so are their colors.
In both cases the lost center on $i=0$ is green, that on $i+j=L$
is blue, and that on $j=0$ is red: substitute the respective side
equation and the qubit residue into the row coordinate.
Two simultaneous side equalities give exactly the corners, and all
three are impossible because $L>0$. This proves the face column.

The primal edges can be checked directly from the cyclic polygons.
A bulk residue-zero qubit has neighbors
\[
 (i-1,j+1),\quad(i+1,j),\quad(i,j-1);
\]
a bulk residue-one qubit has neighbors
\[
 (i+1,j-1),\quad(i-1,j),\quad(i,j+1).
\]
Each pair of incident real faces meets at a different one of these
edges. The missing-color rule gives the three edge colors once each.

On the open red side $j=0$, the residue-zero qubits have $0<i<L$ and
neighbors
\[
 (i-1,1),\quad(i-2,0),\quad(i+1,0).
\]
The first edge is shared by green face $(i-1,0)$ and blue face $(i,1)$
and is red. The second is on the red side of the green face and is
blue. The third is on the red side of the blue face and is green.
The residue-one side qubits have neighbors
\[
 (i,1),\quad(i+2,0),\quad(i-1,0).
\]
Their faces are green $(i+1,0)$ and blue $(i-1,1)$, giving respective
edge colors red, blue, and green again.
The coordinates remain in the domain even next to a corner: the
least positive residue-zero $i$ is $3$, and the greatest residue-one
$i$ is $L-2$.

At $(0,0)$ the sole face is the blue face centered at $(0,1)$.
Its cyclic polygon gives exactly neighbors $(1,0)$ and $(0,2)$.
The edges lie on red and green sides and have colors $g$ and $r$,
respectively. There is no blue edge there.

The rotation $R(i,j)=(L-i-j,i)$ preserves $\Omega_t$, the residue
$i-j\bmod3$, and the cyclic neighbor list up to permutation.
It permutes face colors and side colors by
$g\mapsto r,\ r\mapsto b,\ b\mapsto g$.
It therefore proves the other two open-side cases and other corners.
The same edge lists show that boundary edges join consecutive retained
side qubits, so these cases exhaust the physical boundary.

The red side has $L+1=3t+1$ integer sites, of which $t$ have the
excluded residue $i\equiv2\pmod3$. Thus $|Q_r|=2t+1=d$.
Rotation gives the other counts; the side equations give their intersections.
\end{proof}

\begin{remark}[Corner membership]
A corner belongs to both incident physical sides. Its third,
absent-side color labels its sole real face.
For fixed $c$, the two corners on $\Gamma_c$ and the opposite corner
therefore have different stage-2 incidences.
\end{remark}

\section{The ordinary concatenated-matching decoder}\label{sec:decoder}
\subsection{Two graphs and two bijections}
Fix $c$ and let $\{c_1,c_2\}=\C\setminus\{c\}$.
Vertex unions are tagged disjoint unions. A primal edge used as a
stage-2 vertex is not a stage-2 edge.

\begin{definition}[Stage 1: restricted lattice]\label{def:stage1}
The graph $\LL_{\neg c}^*$ has vertex set
$F_{c_1}\sqcup F_{c_2}\sqcup\{\bd\}$.
For each $e\in E_c$ introduce one labelled edge $\eps_{\neg c}(e)$.
It joins the two real faces separated by $e$, or the one incident real
face to $\bd$ when $e$ is a boundary edge. Thus
\[
 \eps_{\neg c}:E_c\longrightarrow\Delta_1(\LL_{\neg c}^*)
\]
is a bijection.
\end{definition}
This is \cite[Sec.~3, Definition~1]{Lee}; $\neg c$ means color $c$
has been removed, so only the other face colors are real vertices.

\begin{definition}[Stage 2: monochromatic lattice]\label{def:stage2}
The ordinary graph $\LL_c^*$ has vertex set
$F_c\sqcup E_c\sqcup\{\bd\}$.
For each $q\in Q$ introduce exactly one labelled edge $\eps_c(q)$.
If $n_c^F(q)=n_c^E(q)=1$, it joins the unique incident $c$ face to the
unique incident primal $c$ edge. If precisely one object is present,
it joins that object to $\bd$. Lemma~\ref{lem:incidence} ensures
these are all cases. Hence
\[
 \eps_c:Q\longrightarrow\Delta_1(\LL_c^*)
\]
is a bijection. Qubit labels are retained even for parallel edges
with identical endpoint pairs.
\end{definition}
This is \cite[Sec.~3, Definition~2]{Lee}, with its local existence
conditions verified for the chosen family. There is no extra
physical-qubit edge for an exterior triangle.

\subsection{Artificial boundary and matching parity}
For either graph let $\partial_1$ be binary graph incidence and
$\delta$ its restriction to real-vertex rows, deleting the row of $\bd$.
The notation $\MWPM(s;G,\bd)$ means a minimum-weight edge subset $M$
with $\delta M=s$. Every violated real check has odd incidence; all
other real vertices have even incidence. The virtual boundary is
not a measured check, so its parity is not prescribed as part of $s$
\cite[Sec.~2.2, p.~4]{Lee}.
For fixed $s$, the handshake identity determines its incidence parity
as $\sum_v s_v$. Free boundary parity means no extra check equation,
not an independently variable bit.

An edge with only one ordinary endpoint is called \emph{dangling}
before its missing end is attached to $\bd$; afterward it is an
ordinary two-ended graph edge. This describes an incidence deficit,
not an identification of the physical boundary circle with a point.
Different spatial terminations are eligible matching endpoints at
the same artificial node.

\subsection{The two rounds and the physical correction}\label{sec:rounds}
Lee's exact set equations are
\begin{align}
 \widetilde E_{\pred}^{(c)}
 &=\MWPM(\sigma_Z^{(c_1)}\sqcup\sigma_Z^{(c_2)};
                 \LL_{\neg c}^*,\bd),&
 E_{\pred}^{(c)}
 &=\eps_{\neg c}^{-1}(\widetilde E_{\pred}^{(c)}),
 \label{eq:round1}\\
 \widetilde V_{\pred}^{(c)}
 &=\MWPM(\sigma_Z^{(c)}\sqcup E_{\pred}^{(c)};
                 \LL_c^*,\bd),&
 V_{\pred}^{(c)}
 &=\eps_c^{-1}(\widetilde V_{\pred}^{(c)}).
 \label{eq:round2}
\end{align}
Images and inverse images of subsets are setwise bijections.
In the second line $E_{\pred}^{(c)}$ is a set of stage-2 \emph{vertices},
although it was a set of primal \emph{edges} in the first line.
The union combines distinct vertex types.

Write $\delta_{\neg c},\delta_c$ for the two real-row incidence maps,
and $s_c$ for the second-round input. The constraints are
\[
 \delta_{\neg c}\widetilde E_{\pred}^{(c)}
   =\sigma_Z^{(c_1)}+\sigma_Z^{(c_2)},\qquad
 \delta_c\widetilde V_{\pred}^{(c)}=s_c.
\]
For $a=V_{\pred}^{(c)}$, the second equation at a primal-edge vertex
$e=\{u,v\}$ says $a_u+a_v=(E_{\pred}^{(c)})_e$.
At a $c$-face vertex it says $(Ha)_f=(\sigma_Z)_f$.

For each non-$c$ face $f$, its qubits are partitioned into its incident
primal $c$ edges: every qubit of $f$ has exactly one such edge by
Lemma~\ref{lem:incidence}. The local cyclic incidence and the edge-color
rule put that edge in $f$: at a bulk or open-side vertex the only edge
not belonging to $f$ has the color of $f$, and at a corner both edges
belong to its sole face. Distinct $c$ edges cannot share a qubit. Summing their stage-2 equations gives
\[
 (Ha)_f=\sum_{\substack{e\in E_c\\e\subset f}}
                 (E_{\pred}^{(c)})_e=(\sigma_Z)_f
\]
by stage 1. Together with the $c$-face equations this proves
$Ha=\sigma_Z$. This is the real-coordinate version of Lee's
validity argument \cite[Appendix~A.2, Claim~1; A.3, Claim~2]{Lee}.
Validity of two feasible rounds does not imply logical success or
global minimum-weight decoding. The validity argument is conditional on feasible matching outputs; it
is not a separate feasibility or connectivity theorem.

The ordinary algorithm repeats both rounds for $r,g,b$ and returns a
smallest-cardinality candidate among
$V_{\pred}^{(r)},V_{\pred}^{(g)},V_{\pred}^{(b)}$, with a fixed tie rule
\cite[Sec.~3, p.~7]{Lee}. For uniform independent $X$ errors with
$0<p<1/2$, it is most probable among those three candidates.
In the uniform perfect-measurement presentation each round can use
unit edge costs. This is ordinary correction selection, not a
three-color soft-output construction.

\subsection{Chain maps and limits of imported correspondences}
Lee's Appendix~A represents a physical qubit by a dual triangle or
by a hyperedge $h_q$ with one endpoint of each color, completing missing
endpoints with colored exterior labels.
The ordinary dual boundary of a triangle is its three edges; the
hypergraph boundary of $h_q$ is its three vertices. These maps differ.
In the boundary-free source notation, $p_{\mathrm{edge}}^{(c)}h_q$
is the dual $c$ edge between the two other colors and
$p_{\mathrm{vert}}^{(c)}h_q$ is its $c$ vertex. The factorization is
\[
 \partial_1^{\mathcal H}
 =p_{\mathrm{vert}}^{(c)}
  +\partial_1^{\LL_{\neg c}^*}p_{\mathrm{edge}}^{(c)} .
\]
The projectors target the two distinct monochromatic vertex summands
\cite[Appendix~A.1, Eq.~(10)]{Lee}.
With boundaries, the validity equations hold modulo virtual coordinates
\cite[Appendix~A.3, Eq.~(13)]{Lee}.

Delfosse's hypergraph/CSS formulation gives an analogous physical complex
and a projection to restricted surface graphs
\cite[Secs.~4.1--4.3, Definition~5.5, Theorem~5.6]{Delfosse}.
It is not the stage-2 support inverse, and its closed-tiling hypotheses
are not silently extended to corners.
The local Clifford unfolding of a triangular color code gives a folded
surface-code patch with attached layers
\cite[Sec.~III.B, Eq.~(44), Theorem~3]{Kubica}, not an identification
of a layer with $\LL_c^*$.
Meister--Pattison--Preskill distinguish the ordinary matching boundary
from logically distinct surface-code boundary nodes
\cite[Sec.~II.B.1, pp.~2--3]{Meister}; that surface-code logical
characterization is not a theorem about Definition~\ref{def:stage2}.

\section{Boundary structure of the ordinary monochromatic graph}
\label{sec:boundary}
\begin{definition}[Dangling and missing-object classes]\label{def:classes}
For fixed $c$, define the edge subsets
\begin{align*}
 \mathcal B_c^{\dang}
  &=\{e\in\Delta_1(\LL_c^*):\bd\text{ is an endpoint of }e\},\\
 \mathcal B_c^{\missF}
  &=\{\eps_c(q):\mathcal F_c(q)=\varnothing\},\\
 \mathcal B_c^{\missE}
  &=\{\eps_c(q):\mathcal E_c(q)=\varnothing\}.
\end{align*}
The superscripts name the \emph{missing} object: a $\missF$ edge
has a surviving primal-edge vertex and a $\missE$ edge a face vertex.
\end{definition}

\begin{lemma}[Exact dangling criterion]\label{lem:dangling}
The counts $n_c^F(q),n_c^E(q)$ are zero or one, their sum is one or two,
and
\[
 \eps_c(q)\in\mathcal B_c^{\dang}
 \quad\Longleftrightarrow\quad n_c^F(q)+n_c^E(q)=1.
\]
In particular
$\mathcal B_c^{\dang}=\mathcal B_c^{\missF}\sqcup\mathcal B_c^{\missE}$.
\end{lemma}
\begin{proof}
Lemma~\ref{lem:incidence} gives the count bounds. Bulk qubits have both
counts one. At an open side only the face count of its side color
vanishes. At $q_{ab}$ the face counts of $a,b$ vanish and their edge
counts are one; the remaining color has face count one and edge
count zero. Thus $(0,0)$ never occurs. Definition~\ref{def:stage2}
attaches precisely the count-one cases to $\bd$.
\end{proof}

\begin{proposition}[Complete classification]
\label{prop:classification}
For every patch in Definition~\ref{def:coords} and
$\{c,c_1,c_2\}=\C$,
\begin{align}
 \mathcal B_c^{\missF}&=\eps_c(Q_c),&
 |\mathcal B_c^{\missF}|&=d,\label{eq:sideclass}\\
 \mathcal B_c^{\missE}&=\{\eps_c(q_{c_1c_2})\},&
 |\mathcal B_c^{\missE}|&=1.\label{eq:cornerclass}
\end{align}
Consequently
\[
 \eps_c^{-1}(\mathcal B_c^{\dang})
   =Q_c\sqcup\{q_{c_1c_2}\},\qquad
 |\mathcal B_c^{\dang}|=d+1.
\]
\end{proposition}
\begin{proof}
A missing $c$ face occurs exactly on the open $c$ side and at its
two corners $q_{cc_1},q_{cc_2}$, by Lemma~\ref{lem:incidence}.
Their union is $Q_c$.
A missing $c$ edge occurs exactly at the corner whose side colors are
$c_1,c_2$. That corner is not on $\Gamma_c$.
The images are disjoint because $\eps_c$ is a labelled-edge bijection;
their sizes follow from the side count.
\end{proof}

The fixed-color cases are, explicitly,
\begin{center}
\begin{tabular}{@{}lccp{42mm}@{}}
\toprule
Physical location & $n_c^F$ & $n_c^E$ & Stage-2 endpoints\\
\midrule
Bulk & 1 & 1 & $c$ face, primal $c$ edge\\
Open $\Gamma_c$ & 0 & 1 & primal $c$ edge, $\bd$\\
Open $\Gamma_{c_1}$ or $\Gamma_{c_2}$ & 1 & 1 & $c$ face, primal $c$ edge\\
$q_{cc_1}$ or $q_{cc_2}$ & 0 & 1 & primal $c$ edge, $\bd$\\
$q_{c_1c_2}$ & 1 & 0 & $c$ face, $\bd$\\
\bottomrule
\end{tabular}
\end{center}

\begin{remark}[Meaning of the two classes]\label{rem:types}
One class is a whole closed side; the other is its opposite corner.
The opposite corner belongs to both other physical sides. No corner
was assigned arbitrarily to a unique physical boundary.
These are canonical geometric terminal types, not two connected
components of the physical boundary circle or of the merged dangling
star. Their physical logical inequivalence is proved separately in
Section~\ref{sec:logical}.

Lee's source description retains the boundary labels
\[
 \left\{v_{\mathrm{bdry}}^c,\,
        \{v_{\mathrm{bdry}}^{c_1},v_{\mathrm{bdry}}^{c_2}\}\right\}
 \qquad\text{\cite[Appendix~A.3, pp.~22--23]{Lee}.}
\]
A missing $c$ face has the first provenance. At the opposite corner
both non-$c$ faces are missing; the would-be dual edge between them
has the nested-pair provenance. That pair is one virtual stage-2 vertex,
not two physical vertices or an extra data qubit.
Proposition~\ref{prop:classification} identifies the physical origins
of these source labels. The graph construction is given separately
in Section~\ref{sec:resolved}.
\end{remark}

\begin{proposition}[What merging records]\label{prop:merging}
For $S\subseteq\Delta_1(\LL_c^*)$, put
\[
 \tau_c^F(S)=|S\cap\mathcal B_c^{\missF}|\bmod2,\qquad
 \tau_c^E(S)=|S\cap\mathcal B_c^{\missE}|\bmod2 .
\]
The ordinary boundary incidence row records only
\[
 (\partial_1^{\LL_c^*}S)_{\bd}=\tau_c^F(S)+\tau_c^E(S).
\]
It identifies the pairs $(0,0),(1,1)$, and the pairs $(1,0),(0,1)$.
\end{proposition}
\begin{proof}
Every edge at $\bd$ belongs to exactly one missing-object class and
contributes one to that row. Other edges contribute zero.
\end{proof}

Merging removes the separate virtual-endpoint coordinates and the
individual endpoint positions along the physical boundary. Matching
at the common node maintains no separate terminal-type parity constraint.
It need not destroy all provenance: retained qubit labels recover
physical locations, and the tagged type of the surviving endpoint
already distinguishes $\missF$ from $\missE$.
The proposition concerns information in the single boundary incidence
row, not an inability to recover the counts from a labelled edge set.
Neither count is identified with physical logical parity $\lambda_X$.

\section{The boundary-resolved graph}\label{sec:resolved}
\subsection{Meaning of the boundary distinction}
The two nonempty combinatorial classes in
Proposition~\ref{prop:classification} determine the following graph.
Its physical logical interpretation is proved in
Section~\ref{sec:logical}, using the independently known physical
logical operators.

\begin{definition}[Resolved graph and terminal sets]\label{def:resolved}
For fixed $c$, let $G_c=\widetilde{\LL}_c^*$ have tagged vertex set
\[
 \Delta_0(G_c)=F_c\sqcup E_c\sqcup\{b_c^0,b_c^1\},\qquad
 \mathcal B_c=\{b_c^0,b_c^1\},\quad
 \mathcal B_c^{(i)}=\{b_c^i\}.
\]
The boundary objects are two individual virtual vertices. Introduce one
edge $\widetilde\eps_c(q)$ for each physical $q\in Q$, with endpoints
\begin{equation}
 \begin{cases}
 \{f,e\},&
  \mathcal F_c(q)=\{f\},\quad\mathcal E_c(q)=\{e\},\\
 \{b_c^0,e\},&
  \mathcal F_c(q)=\varnothing,\quad\mathcal E_c(q)=\{e\},\\
 \{f,b_c^1\},&
  \mathcal F_c(q)=\{f\},\quad\mathcal E_c(q)=\varnothing .
 \end{cases}                                                   \label{eq:resolved}
\end{equation}
Keep qubit labels, including any parallel edges. These cases exhaust $Q$,
so $\widetilde\eps_c:Q\to\Delta_1(G_c)$ is a bijection.
\end{definition}

All ordinary real vertices and all real edge-endpoint incidences are
unchanged. In particular an ordinary edge with two real endpoints is
unchanged except for its typed resolved-graph label.
The $d$ edges coming from $Q_c$, including its two corners, terminate
at $b_c^0$. The unique edge coming from $q_{c_1c_2}$ terminates at
$b_c^1$. Thus their incident edge sets are disjoint and exhaust the
formerly dangling edges. Neither terminal is a physical qubit or check.
No edge between the two terminals is added.

This realizes Lee's two retained source boundary labels as
$b_c^0=v_{\mathrm{bdry}}^c$ and
$b_c^1=\{v_{\mathrm{bdry}}^{c_1},v_{\mathrm{bdry}}^{c_2}\}$
\cite[Appendix~A.3]{Lee}, using only physical-qubit edges.
The nested source pair is one stage-2 vertex, not an additional edge
of $G_c$. The singleton terminal sets are not two connected components
of the full graph, as the following connectivity result makes explicit.

\begin{definition}[Forgetting the distinction]\label{def:quotient}
The graph map $q_c:G_c\to\LL_c^*$ fixes real vertices and sends both
$b_c^0,b_c^1$ to $\bd$. On edges it sends
$\widetilde\eps_c(q)$ to $\eps_c(q)$.
It is the quotient by identification of the two terminal vertices;
it does not contract a newly introduced edge.
\end{definition}
The endpoint rule \eqref{eq:resolved} and Definition~\ref{def:stage2}
verify this quotient on every edge, including all corners.
Its edge map is bijective and its vertex map is surjective.
Consequently no physical data identity is lost by the quotient's edge map.

\begin{lemma}[Connectivity of the full family]
\label{lem:connected}
For every patch in Definition~\ref{def:coords}, $G_c$ is connected.
In particular its two terminals are joined by a simple graph path.
\end{lemma}
\begin{proof}
It suffices to prove this for red; the rotation of
Lemma~\ref{lem:incidence} gives the other colors.
Red face centers have exactly the form
\[
 f=(i,j)=(1+3m,2+3n),\qquad m,n\geq0,\quad m+n\leq t-1 .
\]
If $n\geq1$, the two qubits $(i,j-1),(i,j-2)$ are the endpoints
of a primal red edge. To check its existence directly from the polygons,
they are consecutive support vertices of both faces centered at
$(i-1,j-1)$ and $(i+1,j-2)$, whose colors are blue and green.
Those centers lie in $\Omega_t$. The first qubit belongs to the red
face $f$, and the second to the red face $(i,j-3)$. Their two resolved
edges therefore give a length-two graph path from $f$ to $(i,j-3)$
through the corresponding primal-edge vertex.

At $n=0$, the analogous primal red edge has endpoints $(i,1),(i,0)$.
It belongs to the blue face $(i-1,1)$ and green face $(i+1,0)$;
these exist since $1\leq i\leq L-2$. The first qubit belongs to
$(i,2)$ and the second lies on the red side. Its two resolved edges
join $(i,2)$ to $b_r^0$. Repeated descent in $n$ therefore connects
every red face vertex to $b_r^0$.

Every primal red-edge vertex has two incident qubit-labelled graph
edges, one for each physical endpoint. Each of their other endpoints
is a red face or $b_r^0$, since the missing-red-edge corner cannot
belong to a primal red edge. Hence every such vertex belongs to the
same component. Finally $b_r^1$ is joined by its corner edge to a
real red face. These are all graph vertices.
For $t=1$ there is just the red face $(1,2)$; the base descent and
corner edge still apply. Finite connected graphs contain a simple
path between any two different vertices, giving the last assertion.
\end{proof}
This is connectivity only. No physical logical class of that path has
been evaluated.

\subsection{Chains, boundary projection, and matching compatibility}
\begin{definition}[Binary incidence and the real-row map]\label{def:chains}
Set $C_i(G_c)=\F^{\Delta_i(G_c)}$ for $i=0,1$ and let
$\partial_c:C_1(G_c)\to C_0(G_c)$ send each edge to the sum of its
two endpoints. No $2$-cells are supplied by this graph notation.
Let
\[
 C_{\mathrm{bdry}}(G_c)=\F^{\mathcal B_c},\qquad
 W_c=\F^{F_c}\oplus\F^{E_c},\qquad
 \rho_c:C_0(G_c)\to W_c,\qquad D_c=\rho_c\partial_c ,
\]
where $\rho_c$ deletes the two terminal coordinates.
Thus $W_c$ is naturally $C_0(G_c)/C_{\mathrm{bdry}}(G_c)$.
Write $q_{c,0}$ and $q_{c,1}$ for the linear maps induced by
Definition~\ref{def:quotient} on vertices and edges.
\end{definition}

\begin{lemma}[Exact incidence compatibility]\label{lem:quotient}
With $\delta_c$ the ordinary real-row map from Section~\ref{sec:rounds},
\begin{equation}
 \partial_1^{\LL_c^*}q_{c,1}=q_{c,0}\partial_c,\qquad
 D_c=\delta_c q_{c,1}.                                      \label{eq:quotient}
\end{equation}
\end{lemma}
\begin{proof}
Both sides of the first equality send a labelled edge to its two
ordinary endpoints, by the three cases of \eqref{eq:resolved}.
Deleting the virtual rows gives the second equality. Extend linearly.
\end{proof}

\begin{proposition}[Matching compatibility with two free terminals]
\label{prop:optimization}
Transport any supplied nonnegative ordinary edge weights by
$\omega_c(\widetilde\eps_c(q))=\omega_c^{\mathrm{ord}}(\eps_c(q))$.
For a fixed real stage-2 syndrome $s_c$, the two optimization problems
\begin{equation}
 \min_{D_c\Gamma=s_c}\ \sum_{e\in\operatorname{supp}\Gamma}\omega_c(e),
 \qquad
 \min_{\delta_c x=s_c}\ \sum_{e\in\operatorname{supp}x}
                                      \omega_c^{\mathrm{ord}}(e)
                                                               \label{eq:opt}
\end{equation}
have bijectively corresponding feasible chains and minimizers under
$x=q_{c,1}\Gamma$, and the same optimum value.
\end{proposition}
\begin{proof}
The edge map is a bijection and preserves every weight.
Equation~\eqref{eq:quotient} identifies their real-row constraints.
Neither problem imposes a terminal parity check. Therefore the
bijection preserves precisely the feasible sets and their costs.
\end{proof}

This proves invariance for these specified problems, not for an arbitrary
change to a matching software interface. If $b_c^1$ is treated as an
ordinary even-parity vertex, or separate prescribed terminal parities
are imposed, the feasible set can shrink. For example at $d=3,c=r$,
the edge labelled $q_{gb}=(0,3)$ joins $b_r^1$ to the sole red face
$(1,2)$. For real syndrome consisting only of that face it is a feasible
unit-cost correction. Requiring even parity at $b_r^1$ rejects it;
the shortest alternative to $b_r^0$ has two edges and cost two.
There is no competing one-edge correction to that face without the
corner edge, since all its other neighbors are real primal-edge vertices.

The operational convention here remains ordinary decoding on $\LL_c^*$,
followed by relabelling its output through $q_{c,1}^{-1}$ for analysis on
$G_c$. A solver supporting two genuinely free boundary vertices could
solve \eqref{eq:opt} directly, but this note changes no decoder.
Equal minimizer sets do not ensure identical tie-breaking, paths used
internally by a solver, shortest-path distances, or dual growth data.
Lee's zero-cost auxiliary boundary connections in a matching reduction
\cite[Sec.~2.2, footnote~1; Appendix~A.3]{Lee} are not physical edges
of $G_c$.

For completeness, Lemma~\ref{lem:connected} implies that $D_c$ is
surjective: for each real vertex choose any graph path to $b_c^0$;
its real boundary is that vertex alone. Binary sums of these paths
realize any $s_c$. This is a graph feasibility result, separately proved
here rather than attributed to Lee's conditional Claims~1--2.

\section{Physical supports and the exact syndrome map}
\label{sec:physicalmap}
\begin{definition}[Physical-support and Pauli maps]\label{def:physicalmap}
Define the linear isomorphism
\[
 T_c:C_1(G_c)\longrightarrow A,\qquad
 T_c\bigl(\widetilde\eps_c(q)\bigr)=q .
\]
For an edge subset or binary chain $\Gamma$, put
\begin{equation}
 V_\Gamma=\operatorname{supp}(T_c\Gamma),\qquad
 \mathscr P_{X,c}(\Gamma)=P_X(T_c\Gamma)=\prod_{q\in V_\Gamma}X_q .
                                                               \label{eq:pauli}
\end{equation}
Here $P_X$ retains its original domain $A$. Writing
$T_c^{\mathrm{ord}}:C_1(\LL_c^*)\to A$ for the linear ordinary
support inverse, the typed relation is $T_c=T_c^{\mathrm{ord}}q_{c,1}$.
Equivalently, the support in \eqref{eq:pauli} is obtained with
$\widetilde\eps_c^{-1}$ applied setwise. No vertex map is inverted.
\end{definition}
Lee already identifies monochromatic edges with physical hyperedges
linearly \cite[Appendix~A.1]{Lee}; this is its explicit physical-qubit
version after the boundary distinction. The inverse is defined on
all $C_1(G_c)$, not only on paths or feasible corrections.

\begin{lemma}[Linearity and multiplication]\label{lem:linearity}
For all $\Gamma_1,\Gamma_2\in C_1(G_c)$,
\[
 T_c(\Gamma_1+\Gamma_2)=T_c\Gamma_1+T_c\Gamma_2,\qquad
 \mathscr P_{X,c}(\Gamma_1+\Gamma_2)
    =\mathscr P_{X,c}(\Gamma_1)\mathscr P_{X,c}(\Gamma_2).
\]
\end{lemma}
\begin{proof}
The first identity follows by extending a basis bijection linearly.
In the second, same-qubit factors cancel because $X_q^2=I$ and
all $X$ factors commute. The equality is exact in this pure sector,
not merely up to phase.
\end{proof}
The same construction with $Z$ is exact as well.
Mixed Pauli multiplication can carry a phase; it is not used here.
CSS separation permits the two syndrome maps to be studied separately,
without a claim of statistical independence.

\begin{definition}[From stage-2 parity data to physical checks]
\label{def:syndromemap}
Let $H_c$ be the $F_c$ row restriction of $H$, and define
$M_c:A\to\F^{E_c}$ by
$(M_ca)_e=a_u+a_v$ for the physical edge $e=\{u,v\}$.
Define $\Lambda_c:W_c\to U$ by
\begin{equation}
 (\Lambda_c(y,z))_f=
 \begin{cases}
 y_f,& f\in F_c,\\
 \displaystyle\sum_{\substack{e\in E_c\\e\subset f}}z_e,
      & f\notin F_c .
 \end{cases}                                                  \label{eq:Lambda}
\end{equation}
All faces in this formula are real. On the second summand this is
the ordinary restricted incidence $\delta_{\neg c}$ after the
linear primal-edge identification $\eps_{\neg c}$, embedded in the
non-$c$ physical face coordinates.
\end{definition}

\begin{lemma}[Graph boundary and physical syndrome; exact on all chains]
\label{lem:syndrome}
For every $\Gamma\in C_1(G_c)$, with $a=T_c\Gamma$,
\begin{align}
 D_c\Gamma&=(H_ca,M_ca),\label{eq:auxiliary}\\
 Ha&=\Lambda_cD_c\Gamma=\Lambda_c\rho_c\partial_c\Gamma .
                                                        \label{eq:syndrome}
\end{align}
In particular the measured $Z$-check syndrome of
$\mathscr P_{X,c}(\Gamma)$ is exactly \eqref{eq:syndrome}.
The full graph boundary also records
\begin{align*}
 \partial_c\Gamma={}&
 \sum_{f\in F_c}(Ha)_f f+\sum_{e\in E_c}(M_ca)_e e\\
 &+\tau_c^F(q_{c,1}\Gamma)b_c^0
   +\tau_c^E(q_{c,1}\Gamma)b_c^1 .
\end{align*}
The last two coordinates are not physical checks.
\end{lemma}
\begin{proof}
At a real face vertex $f\in F_c$, the incident resolved edges are
precisely those labelled by $q\in f$. Their parity is $(Ha)_f$.
At a real primal-edge vertex $e=\{u,v\}$ the incident graph edges
are precisely $\widetilde\eps_c(u),\widetilde\eps_c(v)$,
giving $(M_ca)_e$. This proves \eqref{eq:auxiliary}.

For a non-$c$ face $f$, the local incidence from
Lemma~\ref{lem:incidence} partitions its qubits into primal $c$
edges belonging to $f$, as justified in Section~\ref{sec:rounds}.
Therefore, for arbitrary $a$,
\[
 \sum_{\substack{e\in E_c\\e\subset f}}(M_ca)_e
       =\sum_{q\in f}a_q=(Ha)_f .
\]
Together with the $c$-face row equality, this proves
\eqref{eq:syndrome}. Finally, the rewiring rule identifies the two
terminal rows with the two missing-object parities.
\end{proof}

This is the physical real-row form of the source factorization in
Lee Appendix~A.1, Eq.~(10), and its boundary-coordinate treatment in
Appendix~A.3, Claim~2, Eq.~(13). The local proof makes it valid for
\emph{every} binary chain, not only outputs of a particular decoder.
It also shows why simply writing $HT_c=D_c$ is ill-typed: $D_c$
returns $c$-face checks \emph{and} auxiliary primal-edge parities.
The latter must be passed through restricted incidence to become
physical non-$c$ checks. Neither of the virtual terminals is a
syndrome detector.

If a first-round prediction $z=E_{\pred}^{(c)}$ is compatible with
the non-$c$ measured syndrome and $s_c=(\sigma_Z^{(c)},z)$, then every
resolved chain with $D_c\Gamma=s_c$ gives physical syndrome
$\Lambda_cs_c=\sigma_Z$. Thus the relabelled ordinary correction has
the same physical interpretation, and the stage-1 virtual information
is precisely the $M_c$ constraint on primal-edge vertices.
This virtual information is different from the free terminal rows.

\subsection{Paths, walks, cycles, and general supports}\label{sec:paths}
A simple path is an edge-labelled sequence of distinct vertices with
the chosen connecting edges; as a binary chain it is their sum.
A walk may repeat vertices or edges; its associated chain counts
each traversal modulo two, rather than keeping the set of edges visited
at least once. For example traversing one edge twice gives the zero
chain and identity Pauli. A collection of paths gives their binary sum;
shared edges cancel. A general chain may be disconnected or have
branching support. None of these distinctions restricts the domain of
$T_c$.

Define the endpoint contribution
$\eta_c(v)=\Lambda_c\rho_c(v)\in U$ for a graph vertex $v$.
The three types are
\begin{equation}
 \eta_c(v)=
 \begin{cases}
 f,&v=f\in F_c,\\
 \displaystyle\sum_{\substack{f\notin F_c\\e\subset f}}f,
       &v=e\in E_c,\\
 0,&v\in\mathcal B_c .
 \end{cases}                                                  \label{eq:endpoint}
\end{equation}
Basis faces on the right denote physical syndrome vectors.
An edge vertex can contribute two real non-$c$ checks, or one when
its physical edge is on a boundary.

\begin{proposition}[Endpoint syndrome and the cycle limit]
\label{prop:paths}
If a path or walk has binary chain $\Gamma$ and endpoints $u,v$,
then
\[
 \partial_c\Gamma=u+v,\qquad
 H T_c\Gamma=\eta_c(u)+\eta_c(v).
\]
A cycle chain, or any chain with $\partial_c\Gamma=0$, has zero
physical syndrome. A path between the two different virtual terminals
also has zero physical syndrome.
\end{proposition}
\begin{proof}
In a walk's binary boundary every intermediate occurrence contributes
twice and cancels, leaving its two endpoints (also cancelling if equal).
Apply Lemma~\ref{lem:syndrome}. Both terminal contributions vanish by
\eqref{eq:endpoint}.
\end{proof}
Zero syndrome alone places the support in $Z_X$ without deciding its
class in $Z_X/S_X$. Section~\ref{sec:logical} supplies the additional
physical parity argument that classifies these cycles and terminal paths.
The algebraic cycle condition includes binary unions of cycles, not
just a single simple cycle.

\section{Physical logical topology}\label{sec:logical}
\subsection{Relative chains and physical equivalence}
Continue with the fixed color and pure CSS sector. Real primal-edge
vertices carry the auxiliary stage-1 parity; the terminals have no
measured parity. The exact identities
$D_c=\rho_c\partial_c=\delta_cq_{c,1}$ and $HT_c=\Lambda_cD_c$
are available on every chain.
The physical comparison uses two independent facts:
the triangular patch encodes one qubit
\cite[p.~3, triangular-code construction]{Bombin}, and each complete
physical side supports logical $X$ and $Z$
\cite[Sec.~2.1, p.~4]{Lee}.

\begin{definition}[Physical equivalence and relative difference space]
\label{def:equivalence}
For arbitrary chains $\Gamma_1,\Gamma_2$ define pure-sector equivalence
modulo the physical stabilizer group by
\[
 \mathscr P_{X,c}(\Gamma_1)\sim_{\mathcal S}
 \mathscr P_{X,c}(\Gamma_2)
 \quad\Longleftrightarrow\quad
 T_c(\Gamma_1+\Gamma_2)\in S_X=\operatorname{im}J .
\]
For comparisons within a fixed real stage-2 syndrome, let
\[
 K_c=\ker D_c,\qquad
 R_c=K_c\cap T_c^{-1}(S_X),\qquad \mathcal Q_c=K_c/R_c .
\]
Thus $K_c$ consists of graph chains whose boundary is supported only
on $\mathcal B_c$. By the syndrome lemma $T_cK_c\subseteq Z_X$.
For $z\in K_c$, the graph handshake identity defines
\[
 \beta_c(z)=(\partial_cz)_{b_c^0}=(\partial_cz)_{b_c^1}.
\]
No physical interpretation is built into this definition.
\end{definition}

\begin{lemma}[The physical side detects terminal parity]\label{lem:sideparity}
For every $z\in K_c$ and every nontrivial physical logical-$Z$ support
$\ell_Z$ in this one-qubit code,
\begin{equation}
 \beta_c(z)=Q_c^{\mathsf T}T_cz
          =\ell_Z^{\mathsf T}T_cz=\lambda_c(z).
 \label{eq:target}
\end{equation}
\end{lemma}
\begin{proof}
Exactly the edges labelled by $Q_c$, including its two corners, meet
$b_c^0$. Therefore its incidence coefficient is the sum of the
corresponding qubit coefficients of $T_cz$, namely
$Q_c^{\mathsf T}T_cz$. This uses the proved terminal classification.

The independently known operator $P_Z(Q_c)$ is nontrivial logical $Z$.
There is only one nonzero class in the pure $Z$ logical quotient, so
$\ell_Z+Q_c=Ju$ for some $u\in U$. Since $HT_cz=0$,
\[
 (\ell_Z+Q_c)^{\mathsf T}T_cz
       =u^{\mathsf T}HT_cz=0.
\]
This proves the representative-independent equality.
\end{proof}
The first equality is a direct qubit-incidence calculation; the second
uses the physical code, not a surface-code analogy.

\begin{theorem}[Logical-topology correspondence]\label{thm:logical}
For the stated family, fixed color and pure CSS sector,
\[
 R_c=\ker(\beta_c|_{K_c})=\ker\partial_c .
\]
Every closed resolved graph chain maps to a product of physical face
stabilizers. Every $z\in K_c$ with
$\partial_cz=b_c^0+b_c^1$ maps to the nontrivial physical logical-$X$
class. The map
\[
 \overline T_c:\mathcal Q_c\longrightarrow\mathcal L_X,\qquad
 [z]\longmapsto[T_cz]
\]
is an isomorphism of one-dimensional binary spaces.
In particular all simple paths joining the two terminals are
nontrivial physical logical representatives.
\end{theorem}
\begin{proof}
Pure-sector commutation with a nontrivial logical $Z$ detects the
nonzero $X$ class in a one-qubit code. More explicitly the functional
$\ell_Z^{\mathsf T}a$ vanishes on $S_X$, descends to the
one-dimensional $\mathcal L_X$, and is nonzero because the physical
logical $X,Z$ anticommute. Its kernel on $Z_X$ is exactly $S_X$.
Apply Lemma~\ref{lem:sideparity} to $T_cz\in Z_X$. This gives
$R_c=\ker(\beta_c|_{K_c})$.
For $z\in K_c$, its full boundary is supported on two terminals and
has even total parity. Hence beta zero is equivalent to
$\partial_cz=0$, and beta one is equivalent to boundary $b_c^0+b_c^1$.

The induced map is well-defined, and its kernel is precisely the
defined denominator $R_c$, so it is injective.
Connectivity supplies a simple terminal-joining path $p$.
Its beta is one, and the parity lemma proves its physical image
nonzero. Thus the image contains the unique nonzero class of
$\mathcal L_X$ and the map is surjective.
The same argument applies after exchanging the CSS sectors.
\end{proof}

The theorem covers arbitrary relative chains, not only simple paths.
It depends on the stated one-qubit physical code and the proved terminal
incidences. Connectedness supplies existence of a nonzero representative.
Odd distance and the ordinary boundary assumptions are those of the
selected family; no claim about arbitrary triangular colexes follows.

\begin{corollary}[Converse, path representative, and fixed-fiber equivalence]
\label{cor:representatives}
For $z\in K_c$,
\[
 [T_cz]\ne0\ \text{in }\mathcal L_X
 \quad\Longleftrightarrow\quad [z]\ne0\ \text{in }\mathcal Q_c
 \quad\Longleftrightarrow\quad \beta_c(z)=1.
\]
In that case the support of $z$ contains a simple terminal-joining
path $p$, and $z+p\in R_c$. Every nontrivial physical logical support
$a\in Z_X\setminus S_X$ is stabilizer-equivalent to $T_cp$ for such
a path, but $T_c^{-1}a$ need not itself lie in $K_c$.
For equal real stage-2 syndromes,
\[
 \mathscr P_{X,c}(\Gamma_1)\sim_{\mathcal S}\mathscr P_{X,c}(\Gamma_2)
 \quad\Longleftrightarrow\quad
 \partial_c(\Gamma_1+\Gamma_2)=0 .
\]
\end{corollary}
\begin{proof}
The equivalences follow from the theorem. In the support subgraph of
a beta-one chain the only odd-degree vertices are the two terminals.
Each connected component has an even number of odd-degree vertices;
therefore the two terminals are in the same component. Choose a
simple path in that component. Its removal modulo two leaves a
closed chain, hence an element of $R_c$.
Any other nonzero physical logical support is in the same unique
nonzero class as $T_cp$, so their difference belongs to $S_X$.
The last statement applies the theorem to the fixed-syndrome difference.
\end{proof}

\subsection{Exactly which face products generate the cycles}
\begin{definition}[Physical non-\(c\) face map]\label{def:facemap}
Let $F_{\neg c}=F\setminus F_c$, $U_{\neg c}=\F^{F_{\neg c}}$ and
$\iota_{\neg c}:U_{\neg c}\to U$ be inclusion.
Define the canonical map from real physical faces by
\[
 A_c=T_c^{-1}J\iota_{\neg c}:U_{\neg c}\longrightarrow C_1(G_c).
\]
This definition uses physical stabilizer generators, not an arbitrary
choice of graph cycles.
\end{definition}

\begin{proposition}[Complete cycle basis]\label{prop:facebasis}
The map $A_c$ is injective and
\begin{equation}
 \operatorname{im}A_c=\ker\partial_c=R_c .
 \label{eq:facebasis}
\end{equation}
Thus every closed chain is uniquely the binary sum of the resolved
images of non-$c$ physical faces. No extra corner generators or
relations are required.
\end{proposition}
\begin{proof}
For a non-$c$ face $f$, $HJf=0$. Every primal $c$ edge meets $f$
in both endpoints or neither, by the local incidence used in
Lemma~\ref{lem:syndrome}. Thus $M_cJf=0$ and $D_cA_cf=0$.
Its image is a physical stabilizer, so Theorem~\ref{thm:logical}
implies $A_cf\in\ker\partial_c$.

We now prove equality by a family-wide count and independence check.
The red face coordinates in Lemma~\ref{lem:connected} give
$|F_r|=t(t+1)/2$, and rotation gives the same count for each color.
Counting $\Omega_t$ and subtracting the face centers gives
\[
 n:=|Q|=3t^2+3t+1,\quad
 |F|=\frac{3t(t+1)}2,\quad
 |E_c|=\frac{n-1}{2}.
\]
The last equality counts qubit incidences with primal $c$ edges:
every qubit except the opposite corner belongs to one, and every
such edge has two endpoints.
Because the physical CSS code has $H_X=H_Z=H$ and one encoded qubit,
$1=n-2\operatorname{rank}H$. Hence
$\operatorname{rank}H=(n-1)/2=|F|$, so all real face support vectors
are independent. In particular $A_c$ is injective.

For a finite connected loopless graph, binary incidence has rank
$|\Delta_0|-1$: its image has even coordinate sum, and a spanning
tree attains that rank. Thus
\[
 \dim\ker\partial_c
 =n-(|F_c|+|E_c|+2)+1
 =t(t+1)
 =|F_{\neg c}|=\operatorname{rank}A_c.
\]
The proved inclusion therefore is equality, including $t=1$.
\end{proof}

Under $T_c$ this also gives
$S_X\cap\ker M_c=\operatorname{im}(J\iota_{\neg c})$.
An unrestricted product of physical stabilizers must not be substituted
for this intersection when comparing a fixed stage-2 fiber.
For example $T_c^{-1}Jf$ for a $c$ face is a physical stabilizer chain
but is outside $K_c$: otherwise it would be a combination of non-$c$
faces, contrary to the independence just proved.

\begin{definition}[Induced relative-coordinate complex]\label{def:complex}
The appropriate binary complex for physical equivalence is
\[
 \mathfrak C_c:\quad
 U_{\neg c}\xrightarrow{A_c}C_1(G_c)\xrightarrow{D_c}W_c .
\]
The last group is the graph's vertex space modulo terminal coordinates.
The identities
\[
 D_cA_c=0,\qquad T_cA_c=J\iota_{\neg c},\qquad HT_c=\Lambda_cD_c
\]
make $(\iota_{\neg c},T_c,\Lambda_c)$ a chain map from this complex
to the physical CSS complex $U\xrightarrow J A\xrightarrow H U$.
\end{definition}
Proposition~\ref{prop:facebasis} and Theorem~\ref{thm:logical} prove
\[
 H_1(\mathfrak C_c)=\ker D_c/\operatorname{im}A_c
      =\mathcal Q_c\ \cong\ \mathcal L_X .
\]
This is a precise induced algebraic complex, consistent with the
binary-complex convention in \cite[Definition~3.1]{Delfosse}.
No geometric attachment of $2$-cells is assumed.
By contrast the bare graph pair $(G_c,\mathcal B_c)$ has no $2$-cells;
its relative $H_1$ is $K_c$, of dimension $t(t+1)+1$, and does not
kill the $t(t+1)$ independent physical-stabilizer cycles.
Ordinary homology of the primal disk would also be the wrong object.

\subsection{Cross-checks with physical boundaries, strings, and fibers}
The parity proof used the physical logical $Z$ on the $c$ side as a
detector; it did not assume the inverse of that support was a relative
graph chain. Indeed $T_c^{-1}Q_c\notin K_c$.
Every $q\in Q_c$ belongs to a primal $c$ edge, and no such edge has
both endpoints on $\Gamma_c$ (edges along that side have other colors).
Thus $M_cQ_c\ne0$, even though $HQ_c=0$ and $P_X(Q_c)$ is logical.
This gives a concrete counterexample to dropping the fixed-fiber
condition from the cycle/relative-path characterization.

For either other physical side $\Gamma_a$, $a\ne c$,
$T_c^{-1}Q_a$ is instead a simple terminal-joining path.
To check this combinatorially, order its retained qubits along the
side, as in Lemma~\ref{lem:incidence}.
Its boundary links alternate the two non-$a$ colors.
Adjacent qubits joined by a $c$ link share the corresponding $E_c$
vertex; the alternating link belongs to a $c$ face and the two
qubits share that $F_c$ vertex. Each of these real vertices occurs
once, since a boundary face has exactly two consecutive side qubits.
The first corner on $\Gamma_c$ contributes $b_c^0$; the other corner,
opposite $\Gamma_c$, contributes $b_c^1$.
This verifies the standard boundary logical representatives directly.

More generally a terminal path starts with the opposite-corner qubit
and passes through pairs of qubits at each primal $c$-edge vertex.
Those pairs are the elementary physical $c$ strings of
Section~\ref{sec:strings}; the shared $c$-face excitations cancel as the
path continues to the $c$ boundary. The corner excitation is absorbed
in the same cancellation. Non-$c$ face products deform these relative
paths without changing the fixed auxiliary parities.

The usual logical three-string nets and global physical logical
representatives are equivalent modulo physical face stabilizers
\cite[pp.~2--3, Eqs.~(4)--(6) and Fig.~3]{Bombin}.
The converse corollary therefore assigns their nontrivial physical
class a terminal-path representative. It does not assert that an
arbitrary drawn net is literally a simple path under $T_c^{-1}$.
Deforming a representative into the chosen fiber may use $c$-face
stabilizers, which can change auxiliary parities.
Colored condensation at the physical boundaries remains consistent
with \cite[Secs.~4.2--4.3]{Kesselring}.

The theorem is independent of the numerical stage-1 prediction:
for any fixed real stage-2 syndrome its differences have the same
space $K_c$ and the same quotient. The restriction to that space is
essential. In particular a true physical error with syndrome $\sigma_Z$
need not have $M_c$ parity equal to the decoder's stage-1 prediction.
The result does not claim that stage 1 is correct or that conditioning
on its output produces independent edge errors.

\section{A Meister-type modified decoding graph}
\label{sec:meistergraph}
Only after the physical proof above do we compare with Meister's
modified graph \cite[Sec.~II.B.1, pp.~2--3]{Meister}.
Its relevant structure is a weighted edge-error graph with distinct
virtual boundaries, trivial closed cycles, and nontrivial logical
paths between different boundaries. Here the effective parity map
is $D_c$, not the full physical face matrix $H$.

\begin{proposition}[The exact fixed-branch correspondence]\label{prop:meister}
For the stated patch family and pure CSS sector,
$(G_c,\mathcal B_c,\omega_c)$ is a boundary-resolved stage-2 decoding
graph with the logical structure required of a Meister-type modified
decoding graph, in the following precise sense.
\begin{enumerate}
 \item Its edges biject with physical single-qubit errors through
 $\widetilde\eps_c$, and their Pauli supports are $T_c$.
 \item Its real-vertex constraint is $D_c\Gamma=s_c$; both terminals
 have no prescribed parity. Physical syndrome is $\Lambda_cs_c$.
 \item Relative differences have exactly two physical equivalence
 classes. Closed chains are precisely the fixed-fiber physical
 stabilizer differences, and terminal-joining chains are nontrivial.
 \item Every nontrivial relative class has a simple terminal-joining
 path representative. Identifying the two terminals recovers the
 ordinary graph and its physical edge weights.
\end{enumerate}
\end{proposition}
\begin{proof}
The first two statements are Definitions~\ref{def:resolved},
\ref{def:physicalmap} and Lemma~\ref{lem:syndrome}.
The third and fourth follow from Theorem~\ref{thm:logical},
Corollary~\ref{cor:representatives}, and the explicit quotient
of Definition~\ref{def:quotient}.
\end{proof}

This establishes the same two-terminal relative logical quotient and
cycle equivalence as the planar surface-code construction.
It is not an isomorphism to a particular surface-code lattice or an
equivalence of the full concatenated decoder to a surface-code decoder.
In particular $F_c$ vertices are physical checks while $E_c$ vertices
are stage-1 auxiliary constraints. The reduction of physical string-net
classes to paths occurs in the fixed-color quotient just proved.
Neither Delfosse's projection \cite[Definition~5.5, Theorem~5.6]{Delfosse}
nor Kubica's local-Clifford unfolding
\cite[Sec.~III.B, Theorem~3]{Kubica} is the map used in this proof.

\subsection{Unmodified physical weights}
Use the already adopted transport
\[
 \omega_c(\widetilde\eps_c(q))
       =\omega_c^{\mathrm{ord}}(\eps_c(q))\geq0 .
\]
Splitting changes endpoints, not the physical support or weight of any
edge. There is no physical zero-cost edge joining the terminals.
For uniform independent bit flips with $0<p<1/2$ the usual choice is
$\omega_c=\log((1-p)/p)>0$, or uniform unit cost for ordinary matching.
This is a physical error-support cost, not a summed logical-class
probability or a claim about the distribution conditional on a
predicted stage-1 parity.

\begin{corollary}[Unmodified path-cost correspondence]\label{cor:pathcost}
For supplied nonnegative weights,
\[
 \min_{\substack{z\in K_c\\\beta_c(z)=1}}
       \sum_{e\in\operatorname{supp}z}\omega_c(e)
 =
 \min_{\substack{p\text{ simple path}\\b_c^0\ \mathrm{to}\ b_c^1}}
       \sum_{e\in p}\omega_c(e).
\]
\end{corollary}
\begin{proof}
Every such path is an admissible nontrivial chain.
Conversely Corollary~\ref{cor:representatives} extracts a path contained
in any nontrivial chain. Nonnegativity makes the extracted path cost
no greater. Connectivity and finiteness ensure the sets are nonempty
and the minima exist.
\end{proof}
This is an unmodified shortest-path statement in the fixed-color
relative space, not a minimum over all physical logical supports
with arbitrary weights, and not a cluster-contracted swim distance.

\section{Decoder dual data and resolved clusters}\label{sec:clusters}
Fix the real second-round input $s_c$, and write
$V=\Delta_0(G_c)$, $E=\Delta_1(G_c)$ in this and subsequent metric
sections. These local graph aliases do not denote primal edges.
Initially assume finite strictly positive $\omega_c$.
Let $X_c$ be the interval realization of $G_c$, with shortest-path
metric $d_c$ and edge-length measure $|\cdot|_{\omega_c}$.
The physical edge labels remain available throughout.
Write $\Sigma_c=\operatorname{supp}s_c$ and let $f_c$ be the
returned correction relabelled in $G_c$, so $D_cf_c=s_c$.

\subsection{A free-boundary version of the matching dual}
Meister's clusters use accumulated nonnegative odd-set dual variables,
not connected components of the final correction
\cite[Definition~7, Eq.~(6)]{Meister}. To specify their boundary
meaning, we make the following matching reduction explicit.

\begin{definition}[Boundary-normalized odd-cut dual]\label{def:dual}
For $u,v\in\Sigma_c$, let
\[
 a_u=d_c(u,\mathcal B_c),\qquad
 d_c^\circ(u,v)=\min\{d_c(u,v),a_u+a_v\},\qquad
 \mathcal O_c=\{S\subseteq\Sigma_c:|S|\text{ odd}\}.
\]
For $y=(y_S)_{S\in\mathcal O_c}$ define
\[
 r_u=\sum_{S\ni u}y_S,\qquad
 l_{uv}(y)=\sum_{S:|\{u,v\}\cap S|=1}y_S,\qquad
 Y_c=\sum_S y_S .
\]
The dual is to maximize $Y_c$ subject to
\begin{equation}
 y_S\ge0,\qquad l_{uv}(y)\le d_c^\circ(u,v)\ (u\ne v),\qquad
 r_u\le a_u .
 \label{eq:dual}
\end{equation}
All nonsyndrome vertices, including both terminals, have radius zero.
Certified MWPM data consist of a feasible $f_c$, a feasible $y$, and
the equality $Y_c=\omega_c(f_c)$.
\end{definition}

\begin{proposition}[Exact free-boundary dual and its realization]
\label{prop:dual}
The maximum in \eqref{eq:dual} exists and equals the minimum in
\eqref{eq:opt}. Hence certified data exist for any optimal $f_c$.
The same dual can be formed using ordinary merged-graph distances.
Its radii define balls on $X_c$ without a virtual shortcut between
$b_c^0$ and $b_c^1$.
\end{proposition}
\begin{proof}
Form the complete metric graph on $\Sigma_c\sqcup\{b\}$, where
$d(u,b)=a_u$ and $d(u,v)=d_c^\circ(u,v)$. These are precisely the
distances after identifying the two boundary vertices, hence satisfy
the triangle inequality. Prescribe odd vertices
$T=\Sigma_c$ if $|\Sigma_c|$ is even, and
$T=\Sigma_c\sqcup\{b\}$ otherwise.
A $T$-join means a nonnegative integer edge multiplicity with exactly
these odd-degree vertices.

Its minimum weight equals the physical free-boundary optimum.
Indeed, a physical feasible chain decomposes into edge-disjoint
trails pairing its odd vertices after the terminals are merged.
Delete cycles and replace trails by shortest metric edges; visits to
$b$ may be split there. This gives a metric $T$-join of no greater cost.
Conversely expand metric edges into ordinary shortest paths and add
them modulo two. Real parities are $s_c$ and cancellation cannot
increase cost. Relabel the physical edges in $G_c$. These two
inequalities identify the optima.

The nonnegative odd-cut description of the $T$-join dominant is the
parity-polyhedron result of Edmonds--Johnson
\cite[Sec.~3, pp.~93--95, Eqs.~(3.2), (3.5)--(3.8)]{EdmondsJohnson}.
Its LP dual assigns a nonnegative variable to each $T$-odd cut.
Since $|T|$ is even, complementary shores describe the same odd cut.
Choose the shore excluding $b$ and combine duplicate variables.
Such shores are exactly the odd subsets of $\Sigma_c$.
The edge constraints are then \eqref{eq:dual}: edges $ub$ give
$r_u\le a_u$. Strong duality proves the assertion.
When $\Sigma_c$ is empty the primal and dual values and all radii
are zero, which also covers this case.

Finally a path from $u$ using the merged boundary must first travel
at least $a_u$. Since $r_u\le a_u$, a ball of radius $r_u$ cannot
travel any positive length beyond that boundary via identification.
Consequently its covered portions of physical edge intervals are
exactly those obtained by growing the ball in $X_c$.
It may touch either or both terminals at equality; these are retained
as separate endpoints in $X_c$. Zero-measure terminal points do not
identify the clusters across the two boundaries.
\end{proof}

This is an explicit free-boundary adaptation of Meister's radius
formula. The cut inequalities, rather than unrestricted signed
vertex potentials, justify nonnegative radii. Sparse Blossom's
standard perfect-matching LP permits signed singleton dual variables
\cite[Sec.~2.5.2, p.~8, Eqs.~(4)--(6)]{SparseBlossom}.
The same source explains that its metric-path-graph growth can keep
singleton radii nonnegative (pp.~10--11); signed values are permitted
by the general LP, not asserted to occur in that implementation.
A backend's raw potentials or contracted blossom IDs still
require a verified conversion to \eqref{eq:dual}; clipping negative
potentials is not that conversion. Another mathematically concrete
route is to solve \eqref{eq:dual} as a companion optimization and
check the certificate equality. It need not reproduce the internal
growth history of an earlier PyMatching run. Different optimal duals
may give different clusters, so the dual-selection convention is
part of this soft-output decoder.

\begin{definition}[MWPM cluster set]\label{def:clusters}
For a selected optimal dual in Definition~\ref{def:dual}, put
\[
 \mathcal C_c=\bigcup_{v\in V}B_{r_v}(v)\subset X_c,\qquad
 B_r(v)=\{x:d_c(v,x)\le r\}.
\]
A cluster is a connected component of $\mathcal C_c$.
\end{definition}
This imports the metric-ball convention of Meister
\cite[Definition~1 and Appendix~B]{Meister}, with the preceding
boundary-aware dual supplying its input.
The source's displayed matching primal has singleton equalities,
while its displayed dual assigns all odd sets nonnegative variables.
The nonnegative parity-polyhedron formulation above supplies the
precise LP used here, without silently treating an arbitrary signed
perfect-matching dual as that certificate.
Overlapping balls and balls meeting at a point belong to one cluster.
Radius-zero balls add vertices, not positive-length erased edges.
Pulling back a single connected component from the merged boundary
would incorrectly join components touching different terminals.
The preceding proposition identifies covered edge portions only;
the connected components must be recomputed on the resolved graph.

\subsection{UF and the practical data interface}
For a separately specified UF decoder one can use the same metric
ball definition with the radii from its actual growth run
\cite[Algorithm~1, Definition~8]{Meister}. A resolved-graph convention
is to initialize zero radii at syndrome vertices, grow all syndrome centers
belonging to an odd connected component at equal speed while that
component touches neither terminal, and stop that component when it
becomes even or touches a terminal. Recompute components at collision
events, including simultaneous events. Even components may become
active again after merging with an odd component. Terminal vertices
are absorbing, have radius zero, and never transmit growth to the
other terminal. Stop when every component is even or touches a
terminal, and choose a correction supported on the covered subgraph
by a spanning-tree parity solve. Each terminal-free component then
has even syndrome parity; a terminal-containing component can absorb
its parity, so such a correction exists.
A finite connected positive-length graph ensures termination:
any still-active component expands until it encounters another
syndrome component or a terminal.
This is a defined metric UF variant; it differs from the source's
discrete smallest-odd-cluster schedule. Weighted or discrete backends
must export their actual coverage and growth convention to use it.
It is not asserted to match every UF implementation.
No MWPM optimal-dual bound is asserted for UF.

In the inspected local implementation,
\texttt{ConcatMatchingDecoder.\_decode\_stage2} constructs
\texttt{Matching.from\_check\_matrix} with the supplied check matrix and weights
and calls \texttt{decode\_batch} with \texttt{return\_weights=True}
\cite{Implementation}. The returned prediction and total weight do
not contain $y_S$, $r_v$, or metric coverage. The documented public
decode interfaces expose predictions, matching edges or matched
detectors and weights, without a promised radius/dual certificate
\cite{PyMatching}. Neither final-correction components nor a list
of matched pairs determines the chosen dual.

Required instrumentation is therefore: retain original physical edge
IDs, weights and syndrome-vertex IDs; export a nonnegative odd-cut
certificate satisfying \eqref{eq:dual}, including the original
syndrome membership of each nonzero set; sum its contributions into
$r_v$; and verify feasibility and $Y_c=\omega_c(f_c)$.
A companion dual solver is an alternative, with its cost charged
separately. UF instead exports final radii or covered intervals and
its boundary-stopping convention. This note supplies the mathematical
interface and reference checks; it does not install that instrumentation
in the external circuit-level decoder. A DEM reduction or reweighting
also needs a separate physical-edge and weight audit before it can
instantiate this perfect-measurement theory.

\section{Cluster contraction and the swim distance}\label{sec:swim}
The geometric results below hold for any supplied finite nonnegative
vertex radii on $X_c$, including the specified MWPM and UF radii.
The representative bound in the next section requires the certified
MWPM choice.

\begin{definition}[Contracted length and retained-edge weights]
\label{def:contraction}
Contract each connected component of $\mathcal C_c$ to a point.
Give the resulting interval quotient $\bar X_c$ the length
pseudometric which charges only travel outside $\mathcal C_c$.
For each original labelled edge set
\[
 \bar\omega_c(e)=|e\setminus\mathcal C_c|_{\omega_c},
 \qquad \bar G_c=(V,E,\bar\omega_c).
\]
For a binary chain $z$, $\bar\omega_c(z)$ is the sum of these weights
over its support. The bar does not change $D_c$, $T_c$, or physical
edge identities.
\end{definition}

\begin{lemma}[Partial-edge formula and quotient equivalence]
\label{lem:metric}
Set
\[
 h_v=\max\{0,\max_{u\in V}(r_u-d_c(u,v))\}.
\]
For every edge $e=uv$ of length $w=\omega_c(e)$,
\begin{equation}
 \bar\omega_c(e)=\max\{0,w-h_u-h_v\}.
 \label{eq:partial}
\end{equation}
Distances between original vertices in $\bar X_c$ equal shortest-path
distances in $\bar G_c$.
\end{lemma}
\begin{proof}
A point at coordinate $x\in[0,w]$ on $uv$ can be reached from a vertex
$s$ by entering through $u$ or $v$. Its distance is
$\min\{d_c(s,u)+x,d_c(s,v)+w-x\}$.
Thus the covered part of the edge is the union of its prefix of
length $\min(w,h_u)$ and suffix of length $\min(w,h_v)$.
Their uncovered gap has length \eqref{eq:partial}.
Subdivide every edge at the coverage endpoints. Each covered
subinterval has zero cost and each uncovered one retains its length.
Connected covered components can be contracted without changing
distances: any traversal of a contracted component can be expanded
along a path of zero cost inside that component.
Every new subdivision vertex has degree two before contraction.
Suppressing these vertices replaces successive subinterval costs by
their sum, namely $\bar\omega_c(e)$ on the original edge.
These operations prove the distance identity, including ties and
zero distances.
\end{proof}

In particular an edge is not assigned zero just because its endpoints
belong to the same cluster: that cluster can connect its endpoints
by another route while leaving the middle of this edge uncovered.
Contraction nevertheless permits that other zero-cost route.
An edge of weight $5$ covered for length $1$ at one end and $2$ at
the other retains weight $2$. Boolean zeroing would miscalculate it.

\begin{definition}[Fixed-color stage-2 swim distance]\label{def:swim}
For the selected decoder cluster set define
\begin{equation}
 \phi_c(\mathcal C_c)=
 \operatorname{dist}_{\bar X_c}([b_c^0],[b_c^1])
 =\operatorname{dist}_{\bar G_c}(b_c^0,b_c^1).
 \label{eq:swim}
\end{equation}
Writing $\phi_c(s_c)$ presupposes a fixed decoder and dual/growth
selection rule. The value is a nonnegative real number.
\end{definition}

\begin{theorem}[Exact geometric meaning and physical witness]
\label{thm:swim}
For the stated triangular family,
\begin{equation}
 \phi_c=\min_{\substack{z\in K_c\\\beta_c(z)=1}}\bar\omega_c(z)
       =\min_{\substack{p\text{ simple path}\\ b_c^0\ {\rm to}\ b_c^1}}
                      \bar\omega_c(p).
 \label{eq:swimlogical}
\end{equation}
A minimizing original-edge path $p$ maps under $T_c$ to a nontrivial
physical logical operator. A shortest quotient route can be lifted
to such a physical representative without increasing its cost.
\end{theorem}
\begin{proof}
Lemma~\ref{lem:metric} gives the shortest original-edge path value.
Any relative beta-one chain contains a simple terminal path by
Corollary~\ref{cor:representatives}; removing its other edges cannot
increase nonnegative cost. Conversely each terminal path is such a
chain. This proves \eqref{eq:swimlogical}.
Theorem~\ref{thm:logical} proves $HT_cp=0$ and the nonzero physical
logical class. For a quotient route, expand each contracted traversal
using a path inside the connected cluster, retaining original
terminal endpoints. In the finite subdivided graph this produces
a walk; its binary reduction has boundary $b_c^0+b_c^1$.
Degree-two subdivision vertices have zero boundary, so the reduction
uses either all or none of the segments of each original edge.
It therefore gives an original physical chain. Remove closed
portions to extract a simple terminal path. Its cost cannot increase
and cannot be less than the proved minimum.
\end{proof}

A cluster touching one terminal contracts onto that terminal.
If a connected cluster touches both, their quotient points coincide
and $\phi_c=0$. Its internal path between the \emph{original} terminals
is still a nontrivial physical logical witness. The empty quotient
walk alone is not a physical Pauli operator. Merged clusters are
the connected components of the entire union, not just individual
balls. With positive original lengths, zero swim distance is
equivalent to the two terminals lying in one such component.

Although $G_c$ is connected, no connectedness of its cluster set is
required. In a more general disconnected input, quotient operations
cannot join different components; if the terminal components differ
the extended distance is $+\infty$ and no logical witness is returned.
The present family excludes that case. Negative weights are outside
the metric construction. Original zero weights can instead be
treated as a length pseudometric: retain their labels and incidence,
or contract them for distance computation while retaining expansion
paths and both original terminals. Formula \eqref{eq:partial} still
holds, and the same nonnegative-cost proof applies. Never delete a
zero-cost physical logical chain by forgetting its provenance.

\subsection{Concrete computation and cost}
Given the ordinary labelled stage-2 graph, nonnegative physical
weights, real syndrome, correction and selected decoder radii:
\begin{enumerate}
 \item Restore $G_c$ using Definition~\ref{def:resolved}: edges with
 a missing face meet $b_c^0$, and the opposite-corner edge meets $b_c^1$.
 Relabel the correction through $q_{c,1}^{-1}$ and check $D_cf_c=s_c$.
 No auxiliary boundary-joining edge is carried into the metric.
 \item Obtain radii from the specified dual certificate or UF
 growth run. Record which convention was used. For the analytical
 MWPM claim, require \eqref{eq:dual} and $Y_c=\omega_c(f_c)$.
 Missing growth data are an unavailable soft output, not zero radii.
 \item Compute all $h_v$. Put $R=\max_v r_v$, add a temporary source
 $o$ with arc $o\to v$ of length $R-r_v$, and run Dijkstra on the
 original weighted graph. If its distances are $d_o(v)$, then
 $h_v=\max\{0,R-d_o(v)\}$, since
 $d_o(v)=\min_u(R-r_u+d_c(u,v))$. Remove this temporary source.
 \item Compute $\bar\omega_c$ using \eqref{eq:partial} for each
 labelled edge. This implements interval contraction exactly.
 \item Run Dijkstra on $\bar G_c$ from $b_c^0$ to $b_c^1$.
 Return its distance $\phi_c$, the cluster-data convention, and,
 if requested, a predecessor path with its physical qubit labels.
 Use a predecessor tree or remove zero-cost cycles to obtain a
 simple witness. A zero result is retained.
\end{enumerate}
With radii supplied, two binary-heap shortest-path runs, construction
and linear edge passes take $O((|V|+|E|)\log |V|)$ time and
$O(|V|+|E|)$ storage; a decrease-key Fibonacci heap gives
$O(|E|+|V|\log|V|)$ amortized time. Here $|V|,|E|=O(d^2)$.
The temporary source is an algorithmic device, not a physical edge.

These bounds exclude producing and certifying the decoder data.
For an explicitly stored dual let $L_y=\sum_{S:y_S\ne0}|S|$;
radius accumulation costs $O(L_y)$.
A deliberately simple reference solver computes defect distances in
$O(k(|E|+|V|)\log|V|)$ time and forms the finite LP
\eqref{eq:dual}, with $2^{k-1}$ variables for $k>0$ and $O(k^2)$
constraints. Its explicit representation can cost $O(k^2 2^k)$;
generic rational LP solution is polynomial in that representation
and its bit length, hence exponential in $k$ here. It is a concrete
fallback, not an efficient postprocessing claim.
An instrumented matching solver can avoid this enumeration, but its
certificate conversion, extraction, and verification overhead has
not been implemented or benchmarked here.
For an interval-list UF interface, first merge overlapping intervals
per edge, charging the input size and sorting cost separately.
Near-linear geometric postprocessing alone does not establish
negligible end-to-end overhead.

\section{Certified representative bound and its limits}\label{sec:bound}
Meister's proof uses dual feasibility, odd-cut counting and a
physical nontrivial path, rather than square-lattice distances
\cite[Lemmas~11--12, Theorem~13]{Meister}.
There are two points to supply explicitly here: free boundary
endpoints, and path decompositions that may share vertices.
The source's Definition~4 assumes vertex-disjoint paths; a
minimum-weight relative chain on our graph need not have that form.

\begin{lemma}[Covered weight of any feasible chain]\label{lem:coverage}
For feasible $y$ in \eqref{eq:dual}, form its metric-ball set.
For every binary chain $m$ with $D_cm=s_c$,
\[
 |m\cap\mathcal C_c|_{\omega_c}\ge Y_c .
\]
In particular, certified MWPM data imply
$\bar\omega_c(f_c)=0$.
\end{lemma}
\begin{proof}
At every real vertex, pair its incident half-edges of $m$, leaving
one unpaired exactly when it belongs to $\Sigma_c$. Leave all
terminal half-edges unpaired. Following the pairings partitions
the edges into trails and closed trails; every syndrome vertex
is the endpoint of exactly one open trail. The open trails may
share vertices, but never edges. Terminal-to-terminal trails and
closed trails can be ignored for a lower bound.

Consider a trail with endpoints $u,v\in\Sigma_c$.
Its initial segment of length $r_u$ (or its whole length if shorter)
lies in $B_{r_u}(u)$; its final segment similarly lies in $B_{r_v}(v)$.
Hence covered length along it is at least
$\min\{\text{trail length},r_u+r_v\}\ge l_{uv}(y)$:
its length is at least $d_c(u,v)\ge d_c^\circ(u,v)\ge l_{uv}(y)$,
and $r_u+r_v\ge l_{uv}(y)$ by nonnegativity.
A trail from $u\in\Sigma_c$ to a terminal has length at least
$a_u\ge r_u$ and covered length at least $r_u$.
These estimates count length only once within each trail; vertices
have measure zero, and distinct trails have disjoint edges.

For any odd $S\subseteq\Sigma_c$, the number of these endpoint pairs
with exactly one endpoint in $S$ is odd, hence at least one:
every element of $S$ is an endpoint once, internal pairs contribute
two, and terminal endpoints are outside $S$.
Summing the lower bounds therefore counts each $y_S$ at least once.
It proves the inequality. For $m=f_c$ with
$Y_c=\omega_c(f_c)$, covered weight is at least its entire weight
and cannot exceed it. The uncovered cost is consequently zero.
\end{proof}

\begin{definition}[Two fixed-fiber representative weights]\label{def:gap}
For certified MWPM output $f_c$ define
\[
 W_{c,\mathrm{base}}^{(2)}=\omega_c(f_c),\qquad
 W_{c,\mathrm{opp}}^{(2)}
 =\min_{\substack{D_cm=s_c\\\beta_c(m+f_c)=1}}\omega_c(m).
\]
The opposite set is nonempty by adding a terminal path to $f_c$.
Both optimizations retain the same stage-1 auxiliary parity.
\end{definition}

\begin{theorem}[Meister-type bound for certified stage-2 MWPM]
\label{thm:gap}
For the stated triangular family, fixed color and real stage-2
syndrome, nonnegative finite weights, and the certified dual
cluster construction of Definition~\ref{def:dual},
\begin{equation}
 W_{c,\mathrm{opp}}^{(2)}-W_{c,\mathrm{base}}^{(2)}
 \ \ge\ \phi_c .
 \label{eq:gap}
\end{equation}
For positive weights the clusters use the ordinary interval metric;
for zero weights they use the labelled pseudometric convention above.
No stochastic-noise assumption is required for this cost inequality.
\end{theorem}
\begin{proof}
Let $m$ attain the opposite minimum. Lemma~\ref{lem:coverage} gives
\[
 \omega_c(m)-\omega_c(f_c)
 =|m\cap\mathcal C_c|_{\omega_c}
    -Y_c+\bar\omega_c(m)\ge\bar\omega_c(m).
\]
Since all edges of $f_c$ have zero uncovered cost,
$\bar\omega_c(m)=\bar\omega_c(m+f_c)$.
Their difference belongs to $K_c$ and has beta one.
Theorem~\ref{thm:swim} lower-bounds its uncovered cost by $\phi_c$.
This proves \eqref{eq:gap}.
\end{proof}

The transfer is conditional on an actual exact optimum and the
specified optimal nonnegative dual, but existence of those data for
the mathematical matching problem is proved in Proposition~\ref{prop:dual}.
It is not conditional on an unproved color-code topological analogy.
Ties across logical classes force $\phi_c=0$ under these hypotheses.
Floating-point residuals require an explicit error analysis before
calling a numerical certificate exact.

An arbitrary feasible but suboptimal dual, arbitrary radii, or UF
output does not satisfy this theorem. For example at $d=3,c=r$,
take unit weights and the sole real red-face syndrome.
Its optimal correction is the one-edge corner path of weight $1$;
the opposite class has a two-edge path to $b_r^0$ of weight $2$.
Zero radii give $\phi_r=3$, contradicting a claimed gap bound
$1\ge3$. They are feasible dual data with value zero, but are not
optimal. The certified radius at the excited face is instead $1$,
giving $\phi_r=1$. This isolates the missing hypothesis.

\subsection{What is, and is not, a probability statement}
Equation~\eqref{eq:swimlogical} is an exact geometric minimum over
nontrivial fixed-fiber physical logical chains. As a decoder-confidence
score it is a proxy; calibration is a separate question.
For independent physical bit flips with $0<p_q<1/2$ and
$\omega_c(\widetilde\eps_c(q))=\log((1-p_q)/p_q)$, the ratio of
the probabilities of the two \emph{specific supports} $T_cf_c,T_cm$
is $\exp(\omega_c(m)-\omega_c(f_c))$.
Conditioning on a common event containing both supports cancels its
normalization. This supplies the representative-ratio interpretation
of \eqref{eq:gap}, consistent with Meister's Theorem~13.
It does not assert that a true error belongs to the stage-1 predicted
fiber, nor independence after such a prediction.

Even within the artificial fixed-fiber distribution proportional to
$e^{-\omega_c(m)}$, the class LLR is
\[
 \log\frac{\sum_{D_cm=s_c,\ \beta_c(m+f_c)=0}e^{-\omega_c(m)}}
               {\sum_{D_cm=s_c,\ \beta_c(m+f_c)=1}e^{-\omega_c(m)}} .
\]
These sums include degeneracies, so the representative bound does
not identify or bound this LLR without further argument.
The physical syndrome-conditioned class sum additionally includes
other auxiliary fibers. Meister's exact repetition-code UF identity
\cite[Theorem~10]{Meister} uses the one-dimensional complementary
pair and is not transferred. Neither an upper bound, a constant-factor
approximation, nor an exact posterior failure probability is proved.

The full concatenated decoder also generates stage-1 predictions and
selects among three color branches. Their errors, cross-fiber costs,
branch selection, and any aggregate score are outside \eqref{eq:gap}.
Comparative decoding and general-code postselection already have
distinct definitions in prior work \cite{Postselection,ForcedGap}.
They are not renamed as $\phi_c$ here.
All probability logarithms in this note are natural; a different log
base rescales both physical weights and the swim distance consistently.

\section{Phase-1 conclusions and remaining scope}\label{sec:next}
The construction now runs from physical qubits through the resolved
stage-2 graph, its proved physical logical quotient, actual selected
dual/UF cluster data, and interval contraction to $\phi_c$.
The certified-MWPM representative bound is a separate theorem with
an explicit certificate requirement. Public PyMatching predictions
alone do not implement that requirement.

The physical inputs are Lee's triangular code and decoder maps,
Bombin--Martin-Delgado's logical code structure, and the colored
boundary/string interpretation \cite{Lee,Bombin,Kesselring}.
The projection and unfolding results remain different maps
\cite{Delfosse,Kubica}. The contraction and dual proof strategy adapt
Meister; the free-boundary LP uses the classical parity polyhedron
\cite{Meister,EdmondsJohnson}. The project-specific incidence,
logical-parity, face-basis, lifting and boundary-dual arguments are
proved above, without a priority claim.

Related constructions already extend soft-output methods in other
directions. Kishi et al.\ define bounded and extra-cluster gaps and
prove threshold-restricted comparisons on an already valid logical
decoding graph \cite[Sec.~III.B, Theorems~1--4]{Kishi}.
Lee--English--Bartlett define cluster-statistics confidence metrics
for general QLDPC decoders \cite[Secs.~2--3]{Postselection};
Wills--Yoder--Chuang distinguish summed logical-class likelihoods
from practical forced-decoding weights \cite[Methods]{ForcedGap}.
Smith--Brown--Bartlett also study an exclusive UF surviving-distance rule
\cite[Appendix~E]{Smith}.
Dinc\u{a}--Chan--Benjamin explicitly retain fractional-edge coverage
in their swim-distance definition \cite[Definition~5]{Dinca}.
Chen et al.\ customize PyMatching to extract post-decoding region
weights and compute logical shortest paths for $\mathbb{RP}^2$
cultivation \cite[Sec.~VI, pp.~10--12, Fig.~12]{ChenCultivation}.
That is direct prior work on instrumentation and nonplanar logical
soft output. Their fill-region convention is not identified with our
selected optimal odd-cut certificate, and their antipodal cuts are
not the color-side/corner split of $G_c$.
These works do not supply the physical fixed-color path theorem used
here. This note specializes and justifies an established metric strategy;
it makes no priority claim for cluster contraction, partial coverage,
instrumentation or soft output beyond planar surface codes.
The precise source comparisons and exhaustive claim audit are in the
repository bibliography and review.

Remaining Phase-2 questions are practical dual extraction and its
overhead, effects of stage-1 errors and conditioning, color aggregation,
comparison with the complete decoder's comparative/forced gap,
posterior calibration, and broader noise or code models.
No numerical calibration, aggregate construction, or performance claim is
undertaken in this theory note. The separately authorized circuit-level
extension follows in Part~\ref{part:circuit}.

\input{notes/support/circuit_level_theory.tex}
\input{notes/support/circuit_level_swim_algorithm.tex}

\begin{thebibliography}{99}
\raggedright
\bibitem{Lee}
S.-H. Lee, A. Li, and S. D. Bartlett,
\emph{Color code decoder with improved scaling for correcting circuit-level
noise}, Quantum \textbf{9}, 1609 (2025).
\href{https://doi.org/10.22331/q-2025-01-27-1609}{doi:10.22331/q-2025-01-27-1609}.
Inspected: \href{https://arxiv.org/abs/2404.07482v2}{arXiv:2404.07482v2}.
\bibitem{Bombin}
H. Bomb\'in and M. A. Mart\'in-Delgado,
\emph{Topological Quantum Distillation}, Phys. Rev. Lett.
\textbf{97}, 180501 (2006).
\href{https://doi.org/10.1103/PhysRevLett.97.180501}{doi:10.1103/PhysRevLett.97.180501}.
Inspected:
\href{https://arxiv.org/abs/quant-ph/0605138v3}{arXiv:quant-ph/0605138v3}.
\bibitem{Kesselring}
M. S. Kesselring, F. Pastawski, J. Eisert, and B. J. Brown,
\emph{The boundaries and twist defects of the color code and their
applications to topological quantum computation}, Quantum \textbf{2},
101 (2018).
\href{https://doi.org/10.22331/q-2018-10-19-101}{doi:10.22331/q-2018-10-19-101}.
Inspected: \href{https://arxiv.org/abs/1806.02820v3}{arXiv:1806.02820v3}.
\bibitem{Delfosse}
N. Delfosse, \emph{Decoding color codes by projection onto surface codes},
Phys. Rev. A \textbf{89}, 012317 (2014).
\href{https://doi.org/10.1103/PhysRevA.89.012317}{doi:10.1103/PhysRevA.89.012317}.
Inspected: \href{https://arxiv.org/abs/1308.6207v1}{arXiv:1308.6207v1}.
\bibitem{Kubica}
A. Kubica, B. Yoshida, and F. Pastawski, \emph{Unfolding the color code},
New J. Phys. \textbf{17}, 083026 (2015).
\href{https://doi.org/10.1088/1367-2630/17/8/083026}{doi:10.1088/1367-2630/17/8/083026}.
Inspected: \href{https://arxiv.org/abs/1503.02065v1}{arXiv:1503.02065v1}.
\bibitem{Meister}
N. Meister, C. A. Pattison, and J. Preskill,
\emph{Efficient soft-output decoders for the surface code} (2024).
Inspected: \href{https://arxiv.org/abs/2405.07433v2}{arXiv:2405.07433v2}.
\bibitem{Implementation}
S.-H. Lee, \emph{color-code-stim},
\texttt{src/color\_code\_stim/graph\_builder.py},
methods \texttt{\_build\_triangular\_graph} and
\texttt{\_add\_tanner\_edges}; also
\texttt{src/color\_code\_stim/decoders/concat\_matching\_decoder.py},
\texttt{\_decode\_stage2}, lines 487--541.
\href{https://github.com/seokhyung-lee/color-code-stim/tree/0eb35935c1e5ff30ba3db9def30a9d35bca2f16d}{Commit 0eb35935c1e5},
inspected 11--12 September 2026. Coordinate and interface evidence,
not a mathematical proof.

\bibitem{EdmondsJohnson}
J. Edmonds and E. L. Johnson,
\emph{Matching, Euler tours and the Chinese postman},
Mathematical Programming \textbf{5}, 88--124 (1973).
\href{https://doi.org/10.1007/BF01580113}{doi:10.1007/BF01580113}.
\bibitem{SparseBlossom}
O. Higgott and C. Gidney,
\emph{Sparse Blossom: correcting a million errors per core second
with minimum-weight matching}.
Quantum \textbf{9}, 1600 (2025).
\href{https://doi.org/10.22331/q-2025-01-20-1600}{doi:10.22331/q-2025-01-20-1600}.
Inspected: \href{https://arxiv.org/abs/2303.15933v2}{arXiv:2303.15933v2}.
\bibitem{PyMatching}
PyMatching authors, \emph{Python API Documentation}, stable API
page labelled 2.1.0, inspected 12 September 2026.
\href{https://pymatching.readthedocs.io/en/stable/api.html}{Official API documentation}.
\bibitem{Kishi}
K. Kishi, R. Toshio, J. Fujisaki, H. Oshima, S. Sato, and K. Fujii,
\emph{Even More Efficient Soft-Output Decoding with Extra-Cluster
Growth and Early Stopping} (2026).
Inspected PDF: \href{https://arxiv.org/abs/2602.03336v1}{arXiv:2602.03336v1}.
\bibitem{Postselection}
S.-H. Lee, L. H. English, and S. D. Bartlett,
\emph{Efficient post-selection for general quantum LDPC Codes},
npj Quantum Information \textbf{12}, 96 (2026).
\href{https://doi.org/10.1038/s41534-026-01242-x}{doi:10.1038/s41534-026-01242-x}.
Inspected manuscript: \href{https://arxiv.org/abs/2510.05795v2}{arXiv:2510.05795v2}.
\bibitem{ForcedGap}
A. Wills, T. J. Yoder, and I. Chuang,
\emph{Forced Gap Post-Selection for Quantum LDPC Codes and their Operations}
(2026). Inspected: \href{https://arxiv.org/abs/2605.20346v1}{arXiv:2605.20346v1}.

\bibitem{Smith}
S. C. Smith, B. J. Brown, and S. D. Bartlett,
\emph{Mitigating errors in logical qubits}, Communications Physics
\textbf{7}, 386 (2024).
\href{https://doi.org/10.1038/s42005-024-01883-4}{doi:10.1038/s42005-024-01883-4}.
Inspected: \href{https://arxiv.org/abs/2405.03766v1}{arXiv:2405.03766v1}.
\bibitem{Dinca}
M. Dinc\u{a}, T. Chan, and S. C. Benjamin,
\emph{Error mitigation for logical circuits using decoder confidence}.
Inspected PDF: \href{https://arxiv.org/abs/2512.15689v3}{arXiv:2512.15689v3} (2026).
\bibitem{ChenCultivation}
Z.-H. Chen, M.-C. Chen, C.-Y. Lu, and J.-W. Pan,
\emph{Efficient Magic State Cultivation on $\mathbb{RP}^2$} (2025).
Inspected PDF: \href{https://arxiv.org/abs/2503.18657v1}{arXiv:2503.18657v1}.
\bibitem{StimDEM}
Stim contributors, \emph{Detector error model file format},
\href{https://github.com/quantumlib/Stim/blob/main/doc/file_format_dem_detector_error_model.md}{official documentation},
inspected 12 September 2026, error instructions, separators and offsets.
\bibitem{StimAPI}
Stim contributors, \emph{DetectorErrorModel.shortest\_graphlike\_error},
installed Stim 1.16.0 Python API docstring, inspected 12 September 2026.
\href{https://github.com/quantumlib/Stim/blob/main/src/stim/dem/detector_error_model.pybind.cc}{Official API source}.
\bibitem{CircuitImplementation}
S.-H. Lee, \emph{color-code-stim}, local Phase-2A branch,
commit \texttt{ba6f7dc8b7aaab237d98ad5f08825be4225864c6}.
Inspected \texttt{stim\_utils.py}, \texttt{stim\_symbolic.py},
\texttt{dem\_utils/dem\_manager.py}, \texttt{dem\_utils/dem\_decomp.py},
and \texttt{decoders/concat\_matching\_decoder.py}, 12 September 2026.
Source paths are relative to \texttt{src/color\_code\_stim/};
implementation evidence, not an imported mathematical theorem.
\bibitem{FaultComplex}
T. Hillmann, G. Dauphinais, I. Tzitrin, and M. Vasmer,
\emph{Single-shot and measurement-based quantum error correction via fault complexes},
Phys. Rev. A \textbf{112}, L040401 (2025).
Inspected \href{https://arxiv.org/abs/2410.12963v2}{arXiv:2410.12963v2}.
\bibitem{SpacetimeCodes}
N. Delfosse and A. Paetznick, \emph{Spacetime codes of Clifford circuits}.
Inspected \href{https://arxiv.org/abs/2304.05943v2}{arXiv:2304.05943v2} (2023).
\bibitem{FaultTolerant}
H. Bomb\'in et al., \emph{Fault-tolerant complexes}.
Inspected \href{https://arxiv.org/abs/2308.07844v1}{arXiv:2308.07844v1} (2023).
\bibitem{ColorPipes}
L. S. Herzog, G. Kishony, R. Wille, and A. Fowler,
\emph{Towards Lattice Surgery Compilation for the Color Code Using Pipe Diagrams}.
Inspected \href{https://arxiv.org/abs/2607.05501v1}{arXiv:2607.05501v1} (2026), preprint.
\bibitem{DerksDEM}
P.-J. H. S. Derks, A. Townsend-Teague, A. G. Burchards, and J. Eisert,
\emph{Designing fault-tolerant circuits using detector error models}.
Inspected \href{https://arxiv.org/abs/2407.13826v2}{arXiv:2407.13826v2}.
\bibitem{HomologyCovers}
E. W. Chambers, J. Erickson, K. Fox, and A. Nayyeri,
\emph{Minimum cuts in surface graphs}, author manuscript, Sec.~5.1.
\href{https://ekfox.web.illinois.edu/publications/min-cut-surfaces/min-cut-surfaces.pdf}{Author PDF};
retained local copy inspected 12 September 2026.
\end{thebibliography}
\end{document}
